\documentclass[11pt,a4paper]{article}
\usepackage[T1]{fontenc}
\usepackage[utf8]{inputenc}
\usepackage{lmodern,microtype}
\usepackage[margin=26mm,headheight=15pt]{geometry}
\usepackage{amsmath,amssymb,amsthm,mathtools}
\usepackage{booktabs,longtable,array,enumitem}
\usepackage{xcolor,listings,fancyhdr}
\usepackage[colorlinks=true,linkcolor=blue!45!black,citecolor=blue!45!black,urlcolor=blue!45!black]{hyperref}
\usepackage{bookmark}
\newcommand{\Q}{\mathbb Q}
\newcommand{\C}{\mathbb C}
\newcommand{\Z}{\mathbb Z}
\newcommand{\Gal}{\operatorname{Gal}}
\newcommand{\Tr}{\operatorname{Tr}}
\newcommand{\rad}{\operatorname{rad}}
\newcommand{\Span}{\operatorname{span}}
\newcommand{\id}{\mathrm{id}}
\newcommand{\code}[1]{\texttt{\small #1}}
\DeclareMathOperator{\Res}{Res}
\DeclareMathOperator{\Arg}{Arg}
\newtheorem{theorem}{Theorem}[section]
\newtheorem{proposition}[theorem]{Proposition}
\newtheorem{lemma}[theorem]{Lemma}
\newtheorem{corollary}[theorem]{Corollary}
\theoremstyle{definition}
\newtheorem{definition}[theorem]{Definition}
\newtheorem{remark}[theorem]{Remark}
\lstset{basicstyle=\ttfamily\small,columns=fullflexible,keepspaces=true,
  breaklines=true,breakatwhitespace=false,frame=single,framerule=.3pt,
  rulecolor=\color{black!25},backgroundcolor=\color{black!2},
  showstringspaces=false,aboveskip=10pt,belowskip=10pt}
\setlength{\parskip}{4pt}
\setlength{\parindent}{1.2em}
\setlist{nosep,leftmargin=*}
\pagestyle{fancy}
\fancyhf{}
\fancyhead[L]{\small RadicalRoot: exact algebraic construction}
\fancyhead[R]{\small Reference release 0.1.0}
\fancyfoot[C]{\thepage}
\hypersetup{pdftitle={From Root Objects to Radicals},pdfauthor={OpenAI, prepared for Vladimir Reshetnikov}}
\title{\textbf{From \texttt{Root} Objects to Radicals}\\[5pt]
  \large An exact Galois--Kummer reference solver\\with structural fast paths and explicit branch selection}
\author{Prepared for Vladimir Reshetnikov}
\date{8 September 2026\quad\textbar\quad Reference release 0.1.0}
\begin{document}
\maketitle
\begin{abstract}
An algebraic number represented by a polynomial \texttt{Root} can admit a radical
expression even when \texttt{ToRadicals} leaves it unchanged. The obstruction is
not polynomial degree alone: the exact criterion is solvability of the Galois
group of the minimal polynomial. We turn this criterion into a finite constructive
algorithm. It computes an embedded normal number field, adjoins only the prime
roots of unity required by the group order, builds a prime-cyclic subgroup chain,
and uses trace--Fourier resolvents to construct Kummer generators. Rational linear
algebra expresses every radicand in the preceding field; exact embedded-number
comparisons select principal-power branches. The resulting expression contains
no residual \texttt{Root}, trigonometric, or exponential placeholders.

The accompanying Wolfram Language and Python sources also implement less
expensive structural reductions: power composition, Dickson polynomials,
generalized reciprocal polynomials, pair-sum resolvents, and supplied radical
extensions. The two examples in the motivating question receive especially short
solutions. The normal-field algorithm is mathematically complete with unbounded
exact arithmetic; the delivered software is a bounded reference implementation,
not a claim of practical success on every degree or every backend input. Python
execution tests and Wolfram tests supplied for local execution are distinguished
explicitly.
\end{abstract}
\vspace{4pt}
\noindent\textbf{Release contract.} Exact parameter-free algebraic numbers over
$\Q$; principal complex powers; explicit success, nonsolvability, resource-limit,
and inconclusive outcomes. Radical depth and expression length are not minimized.
\tableofcontents
\newpage

\section{The problem and the deliverable}
The motivating question asks why certain algebraic numbers remain as
\texttt{Root} objects despite possessing radical expressions \cite{question}.
Wolfram's documentation expressly distinguishes its degree-four guarantee from
higher-degree cases that can be missed \cite{toradicals}. A more systematic
solution should therefore do more than apply extra simplification rules.

There are three separate tasks. First, determine whether a radical expression
exists. Second, construct one when it does. Third, ensure that the expression
represents the \emph{selected complex embedding}, not merely an unspecified root
of the same polynomial. This package addresses all three at the algorithmic
level, while exposing resource and backend limitations instead of confusing them
with a mathematical obstruction.

The archive contains a native Wolfram Language package, a separately usable
Python implementation built on SymPy, test suites, examples, an inspection-only
extraction of the supplied notebook's input cells, and this article. The general
backend is implemented in both languages; it is not an instruction to call an
unimplemented ``solve the Galois problem'' routine. It constructs automorphisms
and Kummer generators using explicit field operations. The Python implementation
has been executed in the preparation environment. The Wolfram evaluation service
returned HTTP 404, so its package has not been executed in a Wolfram kernel here.
The supplied \code{.wlt} tests make that missing validation reproducible locally.

\subsection{What ``any Root'' means here}
An input is a fixed $\alpha\in\overline{\Q}\subset\C$, represented with an exact
embedding. The Wolfram interface accepts an exact algebraic expression, including
a polynomial \texttt{Root}; the Python interface accepts a univariate rational
polynomial and a SymPy root index. Parameters, floating-point coefficients,
transcendental input, and inverse-function \texttt{Root} objects that are not
algebraic numbers over $\Q$ are outside this release's contract.

A reducible defining polynomial must first be reduced to the irreducible factor
containing the selected root. For example, a rational root of
$(x-7)(x^5-x-1)$ is certainly radical-expressible even though the other factor has
a nonsolvable Galois group. The Python interface uses the selected
\texttt{CRootOf}'s factor; the Wolfram interface computes a minimal polynomial.

\subsection{Output syntax and branch semantics}
\begin{definition}
A radical expression is generated by rational constants, $i$, addition,
multiplication, and rational powers, with all subexpressions well-defined.
Negative integer powers include division. Each noninteger power is the principal
complex power. The implementation's grammar test rejects algebraic-root objects,
free symbols, floating-point literals, and trigonometric or exponential
functions.
\end{definition}
Thus $\cos(2\pi/7)$ is not accepted as an answer merely because it is algebraic.
A primitive root of unity, however, has the allowed representation
\begin{equation}\label{eq:zeta}
 z_n=(-1)^{2/n}\qquad(n>1),\qquad z_1=1.
\end{equation}
Since the principal argument of $-1$ is $\pi$, this is precisely
$\exp(2\pi i/n)$. Equation~\eqref{eq:zeta} is used as a radical, not as an
unevaluated exponential. It may not be a shortest radical expression for the
same constant. Real targets are allowed to have complex intermediate values;
requiring every intermediate radical to remain real is a different problem.

\section{What the supplied notebook contributes}
The supplied archive contains one notebook, \code{FindExtension.nb}, of 986,737
bytes. A non-evaluating box reader extracted 220 input cells for inspection.
The extraction is deliberately not advertised as executable source: certain
special boxes and incomplete scratch cells retain display-oriented syntax.
The original notebook was read, not executed. Source locations below refer to
line numbers in its textual notebook representation, not input-cell numbers.

The experiments contain two valuable ideas: search for coefficients of low-degree
factors, and use combinations of conjugates to discover smaller extension
fields. Those are recognizable fragments of resolvent and fixed-field methods.
A systematic package should retain these ideas while replacing unbounded search,
implicit branch choices, and state-dependent notebook execution.

\begin{longtable}{@{}p{.22\textwidth}p{.69\textwidth}@{}}
\toprule
Notebook component & Interpretation and release response\\
\midrule
\endhead
\code{findExtension}, line 65 & Searches for a coefficient of a quadratic
factor by elimination. A pair-sum resolvent gives a direct exact polynomial for
such coefficients. A small coefficient field is a useful candidate, not by
itself a proof that the chosen target has been solved.\\[4pt]
\code{findExtension3}, line 126 & Analogous cubic-factor experiments suggest
higher subset resolvents. The release implements pair sums only; arbitrary
subset resolvents are discussed as a future fast-path improvement.\\[4pt]
\code{project}, line 182 & Real and imaginary parts of conjugates may have
smaller minimal polynomials. A successful projection still needs a target
reconstruction step; failure says nothing about solvability.\\[4pt]
\code{factor}, line 252; \code{solve}, line 315 & Candidate selection uses
\code{PossibleZeroQ}, or the first real solution. The release uses exact
embedded-number equality and explicit failure outcomes. A ``first real'' rule
cannot generally preserve an arbitrary input root.\\[4pt]
\code{fuzz}, line 436 & Random subsets, 400-digit approximation and
\code{RootApproximant} are useful exploratory tools. The displayed infinite loop
has no exhaustion exit even after its finite subset space is exhausted. No
finite failure criterion or complete search schedule is supplied. The general
solver instead uses a finite trace-basis search.\\[4pt]
\code{radicalDepth}, line 21 & This is a syntactic score, not a solvability test.
In particular a residual \code{Root} can fall into its depth-zero case. The new
grammar test excludes residual roots before accepting an answer.\\[4pt]
\code{simplify}, line 560 & Randomized transformations, an expanding time budget,
and private simplification scores are separated from construction in the release.
The earlier selectors spell \code{Simplify`SimpifyCount}, whereas this cell uses
\code{Simplify`SimplifyCount}; the misspelling should not be relied on as a
complexity metric. Neither private name is used in the package.\\
\bottomrule
\end{longtable}

For a monic quadratic divisor $x^2+px+q$, remainder vanishing is equivalently
vanishing of its two coefficient expressions. This is simpler to specify and
audit than negating an elimination of a nonvanishing remainder. Moreover,
$p=-(r_i+r_j)$ for two roots. A repeated pair sum can correspond to different
unordered pairs, so even knowing $p$ need not determine $q$ uniquely. Such
collisions must not be erased by careless square-free processing.

There are also practical notebook concerns: unrestricted use of global $x$,
empty selections followed by \code{[[1]]}, cells depending on earlier interactive
state, and expression-size growth. These are reasons to use explicit arguments,
lexically scoped variables, and structured result records. They do not invalidate
the successful identities found experimentally.

The notebook's large degree-nine, degree-eighteen, and degree-thirty experiments
were not all converted into regression fixtures or benchmarked here. The release
includes one explicit degree-eight example from it in addition to the motivating
question. No claim of complete notebook-corpus coverage is made.

\section{Solvability and the normal closure}
Let $f\in\Q[x]$ be the monic irreducible minimal polynomial of $\alpha$, let $L$
be its splitting field inside our fixed copy of $\C$, and put
\[
 G=\Gal(L/\Q),\qquad d_L=[L:\Q]=|G|.
\]
The classical solvability criterion and the cyclic-extension theorem are the
starting points, not new results of this release \cite[Prop.~5.27,
Thm.~5.34]{milne}. The construction below spells out the finite searches and
coordinate operations needed to turn them into code.

\begin{theorem}\label{thm:criterion}
The selected $\alpha$ has a radical expression over $\Q$ if and only if $G$ is
solvable. For an irreducible $f$, either every conjugate is radical-expressible or
none is.
\end{theorem}
\begin{proof}
A radical tower containing $\alpha$ can be enlarged by roots of unity and the
conjugates of its generators so that its normal closure has a solvable Galois
group. One can see this by adjoining the conjugates layer by layer: after the
preceding layer and the needed roots of unity have been included, the new layer
is generated by roots of binomials and contributes an abelian kernel. The
splitting field of $f$ is a normal subfield of this enlarged field, so its Galois
group is a quotient of a solvable group. This proves necessity for one selected
root, not merely for a prior promise about all roots.

Conversely, Sections~\ref{sec:cyclotomic}--\ref{sec:tower} construct a radical
tower containing $L$ when $G$ is solvable and give an expression for the specific
embedded $\alpha$. Applying the same final coordinate step to any conjugate
gives the last assertion.
\end{proof}

An extension $\Q(\alpha)/\Q$ need not be normal even when $\alpha$ is easily
expressible by radicals. Its automorphism group alone is therefore not the
correct obstruction. For instance, $\Q(\sqrt[3]{2})$ has too few internal
automorphisms to encode the full cubic splitting problem. The normal closure
is essential to the completeness argument.

\section{An exact representation of the normal field}
All roots of $f$ can be represented as exact algebraic numbers without first
finding radicals. A primitive-element operation on the list gives
$L=\Q(\theta_L)$. This is not circular: exact algebraic-root isolation, field
membership, and factoring in number fields are permitted backend operations;
radical extraction for arbitrary algebraic numbers is not assumed.

For an embedded field $M=\Q(\theta)$ with minimal polynomial $P$ of degree $D$,
represent an element by its unique polynomial of degree less than $D$ in
$\Q[X]/(P)$. Addition and multiplication are rational polynomial operations
followed by reduction modulo $P$. An inverse is computed by the Euclidean
algorithm. Equality in this representation is exact equality of coefficient
vectors. The fixed complex embedding is carried by the choice of $\theta$. The CAS
representations used here are documented in \cite{wlexact,sympydomains}.

\subsection{Computing automorphisms without a Galois-group oracle}
If $M/\Q$ is normal, factor $P$ over $M$. It must have $D$ distinct linear
factors. Write their roots as $t_g(\theta)$ with $\deg t_g<D$. Then
\begin{equation}\label{eq:action}
 g\bigl(h(\theta)\bigr)=h\bigl(t_g(\theta)\bigr),
 \qquad h\in\Q[X]/(P).
\end{equation}
Every embedding is an automorphism, and every automorphism appears once. A
factorization failing the linearity or root-count check is a backend failure,
not a reason to silently continue.

Composition deserves care. With the convention $(g\circ h)(a)=g(h(a))$, the
image polynomial for $g\circ h$ is
\[
 t_h(t_g(X))\pmod {P(X)}.
\]
The apparent reversed order is intentional. The group table uses exactly this
convention. Explicit substitution makes the table associative; the Python
general-purpose finite-group test helper also checks associativity when given
an external table.

\subsection{A Wolfram-specific primitive-element requirement}
The call \code{ToNumberField[list, All]} is used, not its default second
argument. The documentation says that \code{All} chooses the smallest common
extension, whereas the automatic form need not do so \cite{tonumberfield}.
An unrelated larger field could invalidate the intended normal-field or degree
arguments and needlessly inflate the computation.

The package also makes the chosen generator integral. If
$a_nX^n+\cdots+a_0$ is its primitive integer minimal polynomial, then
$a_n\theta$ is integral and generates the same field. The documentation promises
representation in the specified generator when that generator is integral.
The encoder checks that it has not silently changed the generator before using
\code{AlgebraicNumberPolynomial}. This avoids interpreting coordinate
coefficients relative to the wrong primitive element.

\section{A small sufficient cyclotomic base}\label{sec:cyclotomic}
Set
\begin{equation}\label{eq:fields}
 m=\rad(d_L)=\prod_{p\mid d_L}p,\qquad
 K=\Q(z_m),\qquad M=LK.
\end{equation}
Only distinct prime divisors are needed: the subgroup chain below has
\emph{prime} cyclic quotients, not arbitrary prime-power cyclic quotients.
The value $m=1$ is used for $d_L=1$. This is a sufficient conductor, not a claim
of a minimal conductor. For example the factor $2$ can sometimes be omitted
without changing the cyclotomic field.

Both $L$ and $K$ are normal over $\Q$, hence so is $M$. Let
\[
 H=\Gal(M/K),\qquad D=[M:\Q],\qquad d_0=[K:\Q]=\varphi(m).
\]
Starting with all automorphisms from~\eqref{eq:action}, retain exactly those
fixing the stored field element $z_m$. This gives $H$, and the implementation
checks
\begin{equation}\label{eq:stabilizer}
 |H|\varphi(m)=D.
\end{equation}

\begin{proposition}\label{prop:relative}
The groups $H$ and $G$ are solvable simultaneously. Every prime dividing $|H|$
divides $m$.
\end{proposition}
\begin{proof}
Restriction identifies $H$ with $\Gal(L/L\cap K)$, a normal subgroup of $G$.
The field $L\cap K$ is a Galois subfield of the abelian extension $K/\Q$;
therefore $G/\Gal(L/L\cap K)$ is abelian. A subgroup of a solvable group is
solvable, and an extension of a solvable group by an abelian group is solvable.
Thus the equivalence follows. Finally $|H|$ divides $|G|=d_L$, so every prime
factor of $|H|$ occurs in the product defining $m$.
\end{proof}

This observation prevents a common incorrect shortcut: adjoining roots of
unity cannot turn a genuinely nonsolvable polynomial over $\Q$ into a solvable
one while hiding the obstruction in an unexamined quotient. Here the quotient
is explicitly known to be abelian.

\section{The finite subgroup calculation}
Compute the derived series
\[
 H^{(0)}=H,\qquad H^{(j+1)}=[H^{(j)},H^{(j)}].
\]
Because the groups are finite, the orders either decrease or stabilize. A
nontrivial stable derived subgroup certifies nonsolvability. This is the only
kind of conclusion called \code{NotSolvable} by the generic group path.

For a solvable current subgroup $S$, construct a maximal proper subgroup
containing its commutator subgroup as follows. Start with $N=[S,S]$. Scan the
elements $g\in S$ in fixed order. If $\langle N,g\rangle$ is still proper in
$S$, replace $N$ by that subgroup; otherwise retain $N$. The scan needs no
backtracking. Once $\langle N,g\rangle=S$, increasing $N$ cannot make the same
join proper again.

\begin{lemma}
At the end of the scan, $N\triangleleft S$ and $S/N$ is cyclic of prime order.
\end{lemma}
\begin{proof}
Every intermediate subgroup contains $[S,S]$ and is therefore the inverse
image of a subgroup of the abelian group $S/[S,S]$. In particular it is normal
in $S$. The preceding monotonicity observation proves maximality after one
scan. The quotient $S/N$ is a nontrivial simple abelian group, hence cyclic of
prime order.
\end{proof}

Repeat with $S=N$ until the identity subgroup is reached. This produces
\begin{equation}\label{eq:chain}
 H=H_0\triangleright H_1\triangleright\cdots\triangleright H_r=\{1\},
 \qquad [H_{i-1}:H_i]=p_i\text{ prime}.
\end{equation}
Each subgroup only has to be normal in the immediately preceding subgroup,
not in all of $H$. Explicit normality and prime-index checks are retained in
the code even though the construction proves them. The corresponding fixed
fields satisfy
\[
 K=M^{H_0}\subset M^{H_1}\subset\cdots\subset M^{H_r}=M,
\]
with each adjacent extension cyclic of degree $p_i$.

\section{A deterministic Kummer generator}
Fix one step $S=H_{i-1}$, $N=H_i$, $p=[S:N]$, and choose
$\sigma\in S\setminus N$. Its restriction generates $\Gal(M^N/M^S)$.
Let $\zeta=z_m^{m/p}$, a primitive $p$th root of unity lying in $K$.
For $k=0,\ldots,D-1$, form
\begin{align}
 a_k&=\sum_{\nu\in N}\nu(\theta^k),\label{eq:trace}\\
 \rho_k&=\sum_{j=0}^{p-1}\zeta^{-j}\sigma^j(a_k).\label{eq:fourier}
\end{align}
Choose the first nonzero $\rho_k$ by exact field arithmetic. No random
coefficient bounds or guessed algebraic approximants are involved.

\begin{lemma}[Finite nonvanishing search]\label{lem:nonzero}
At least one of the $D$ values in~\eqref{eq:fourier} is nonzero.
\end{lemma}
\begin{proof}
The elements $1,\theta,\ldots,\theta^{D-1}$ span $M$ over $\Q$. The map
$T_N:b\mapsto\sum_{\nu\in N}\nu(b)$ is a $\Q$-linear surjection onto $M^N$:
for $b\in M^N$, $T_N(b/|N|)=b$. Consequently the $a_k$ span $M^N$ over $\Q$.

The distinct automorphisms $1,\sigma,\ldots,\sigma^{p-1}$ of $M^N/M^S$ are
linearly independent as maps over $M^N$ (independence of field homomorphisms).
Therefore the map $A=\sum_j\zeta^{-j}\sigma^j$ is not the zero map. If every
$A(a_k)$ vanished, $\Q$-linearity and the spanning property would force $A=0$,
a contradiction.
\end{proof}

\begin{proposition}[Constructive cyclic radical step]\label{prop:kummer}
For the chosen $\rho=\rho_k\ne0$,
\[
 \nu(\rho)=\rho\ (\nu\in N),\qquad
 \sigma(\rho)=\zeta\rho,\qquad
 c=\rho^p\in M^S,
\]
and $M^N=M^S(\rho)$.
\end{proposition}
\begin{proof}
Normality of $N$ in $S$ implies that $\sigma$ preserves $M^N$, and each $a_k$
belongs to $M^N$. Thus so does $\rho$. Reindexing the Fourier sum gives
$\sigma(\rho)=\zeta\rho$, using $\sigma^p(a_k)=a_k$ and $\sigma(\zeta)=\zeta$.
Taking $p$th powers shows that $c$ is fixed by $\sigma$ and by $N$, hence by $S$.
Since $\zeta\ne1$ and $\rho\ne0$, $\rho$ is not fixed by $\sigma$ and does not
belong to $M^S$. The intermediate degree divides the prime $p$ and is not one;
it is therefore $p$, proving the field-generation claim.
\end{proof}

The code rechecks these eigenvector and fixed-field identities modulo $P$.
This finite basis search is the principal replacement for the notebook's
unbounded randomized searches. It does not guarantee a small generator; it
guarantees a generator once the exact field and group data are available.

\section{From abstract roots to principal powers}
A relation $\rho^p=c$ by itself is insufficient for an explicit answer using
principal powers. It identifies $p$ choices, not the chosen embedded $\rho$.
Let $q=c^{1/p}$ be the principal root of the exact embedded value $c$. Since
$q^p=\rho^p$ and $c\ne0$,
\[
 \rho=\zeta^b q
\]
for exactly one $b\in\{0,\ldots,p-1\}$. All these candidates lie in $M$.
The implementation finds $b$ using exact embedded field membership and equality.
It does not accept a small floating-point residual as proof.

The Wolfram version computes \code{RootReduce} on each difference and requires
literal zero. The Python version re-embeds the principal power in its exact
algebraic field and compares the reduced coordinate arrays. These checks trust
the respective CAS number-field and embedding machinery. They are not a
formal proof object verified independently of that machinery. Arbitrary
precision may be used internally by a CAS to identify embeddings, but no
user-level tolerance test is substituted for an algebraic equality.

\begin{remark}[A useful failure test]
The two numbers $\sqrt2$ and $-\sqrt2$ have the same minimal polynomial.
Consequently checking only $f(e)=0$, or equality of minimal polynomials,
would accept a wrong conjugate. Both test suites explicitly reject this case.
\end{remark}

No \code{PowerExpand} operation is used. In particular, the generally false
principal-power identity $(uv)^{1/p}=u^{1/p}v^{1/p}$ is never needed. The only
branch adjustment is the finite verified multiplier $\zeta^b$.

\section{Coordinates, a radical tower, and output}\label{sec:tower}
Maintain two parallel objects: a rational basis $B_i$ of $M^{H_i}$ as exact
elements of $M$, and a symbolic basis $\mathcal B_i$ giving their already
constructed radical expressions. Initially
\[
 B_0=(1,z_m,\ldots,z_m^{d_0-1}),\qquad
 \mathcal B_0=(1,z_0,\ldots,z_0^{d_0-1}),
\]
where the first straight-line definition is $z_0=(-1)^{2/m}$ (or $z_0=1$).

At the $i$th step, write the exact element $c_i=\rho_i^{p_i}$ in the preceding
basis:
\begin{equation}\label{eq:coords}
 c_i=\sum_{b\in B_{i-1}}q_{i,b}b,\qquad q_{i,b}\in\Q.
\end{equation}
The system is overdetermined in ambient coordinates but has a unique solution:
the basis is independent and $c_i$ belongs to the lower field. A rational
linear solver computes these coefficients; the implementation checks the
reconstructed vector exactly. Define
\begin{equation}\label{eq:definition}
 r_i=z_0^{b_i m/p_i}
 \left(\sum_{b\in B_{i-1}}q_{i,b}\,\mathcal B_{i-1}(b)\right)^{1/p_i}.
\end{equation}
By the branch check, its embedded value is exactly $\rho_i$.
Update the bases to the ordered lists
\begin{equation}\label{eq:basisupdate}
 B_i=(b\rho_i^j)_{0\le j<p_i,\ b\in B_{i-1}},\qquad
 \mathcal B_i=(\mathcal B_{i-1}(b)r_i^j)_{0\le j<p_i,\ b\in B_{i-1}}.
\end{equation}
The product-basis theorem and Proposition~\ref{prop:kummer} prove independence.
The size check is
\[
 |B_i|=\varphi(m)\prod_{j=1}^i p_j=D/|H_i|.
\]
At the end $B_r$ is a rational basis of $M$. Solve one more coordinate system
for $\alpha$, then substitute the definitions of $z_0,r_1,\ldots,r_r$ in order.

\begin{theorem}[Constructive correctness and termination]\label{thm:main}
Given terminating exact algorithms for rational polynomial factorization,
embedded algebraic-number arithmetic, primitive elements, and factorization
over number fields, the unbounded normal-field procedure terminates. It either
proves nonsolvability or returns a radical expression equal to the selected
$\alpha$. It returns a radical expression precisely when one exists.
\end{theorem}
\begin{proof}
The initial splitting field and cyclotomic compositum have finite degree, and
the number-field operations have the stated terminating contracts. The group
and each subgroup scan are finite. If the derived series stabilizes
nontrivially, Proposition~\ref{prop:relative} and Theorem~\ref{thm:criterion}
prove nonsolvability. Otherwise all subgroup indices are prime and divide $m$.
Lemma~\ref{lem:nonzero} makes each resolvent search finite, and
Proposition~\ref{prop:kummer} establishes the required cyclic radical steps.

Induction on~\eqref{eq:basisupdate}, using the exact branch equality at each
step, proves that every symbolic basis expression has the value of its paired
field element. Equation~\eqref{eq:coords} always has the claimed unique rational
solution. At the final step the coordinate expression therefore equals
$\alpha$. Its definitions use only allowed arithmetic and rational powers,
with nonzero Kummer radicands. This proves correctness and existence of all
required steps. The solvability criterion proves the converse.
\end{proof}

The conditional backend clause is important. A mathematical decision algorithm
is not a guarantee that a particular version of a CAS will successfully finish
every call before memory exhaustion, a time limit, or an implementation
exception. The delivered wrappers expose those outcomes. Their default bounded
settings are \emph{not} a complete decision procedure in finite practical time.

\subsection{Straight-line output versus full expansion}
A nested radical can contain the same large subexpression repeatedly. The
result records retain symbolic definitions and a final value expression so
that a consumer can inspect a straight-line program instead of reading one
enormous expanded formula. The normal-field implementations also construct
an expression with all temporary symbols substituted. This can itself be
expensive and there is no expression-size minimizer in this release.

The certificates contain chain indices, branches, radicands, eigenvector
coordinates and a group multiplication table. They are audit records, not a
standalone portable formal certificate format: an independently trusted
verifier would additionally need a complete serialized embedded-field
presentation and embedding certificates. Python's \code{verify\_result}
provides an independent final CAS equality check; the Wolfram routine always
performs its final \code{RootReduce} equality check before reporting success.

\section{Structural fast paths}
A normal closure is often unnecessary for a short answer. The solver first
tries exact structural transformations whose candidate results are verified
against the selected root. Exhausting them yields ``not found,'' never
``not solvable.''

\subsection{Low degree and affine power composition}
The usual degree-at-most-four formulas supply initial candidates. A backend
may leave a formula incomplete or retain a root object; the grammar and equality
checks prevent such an object from being counted as successful radical output.

For monic $f$ of degree $n$, let $h=-[x^{n-1}]f/n$, so that $f(x+h)$ is
centered. Take the greatest common divisor $d$ of the positive exponents with
nonzero coefficients. If $d>1$, then
\[
 f(x+h)=g(x^d),\qquad \deg g=n/d.
\]
Solve $g(y)=0$, and for each candidate $\beta$ try
\[
 h+z_d^k\beta^{1/d},\qquad 0\le k<d.
\]
Recursion strictly reduces polynomial degree. The fast-path recursion has an
additional safety depth limit in this release, so it remains a bounded
candidate generator, not the source of the completeness theorem.

\subsection{Dickson polynomials and correlated radicals}
Define
\[
 D_0(x,a)=2,\quad D_1(x,a)=x,\quad
 D_n(x,a)=xD_{n-1}(x,a)-aD_{n-2}(x,a).
\]
Induction gives the central identity
\begin{equation}\label{eq:dickson}
 D_n(u+a/u,a)=u^n+(a/u)^n\qquad(u\ne0).
\end{equation}
If a centered polynomial is exactly $D_n(x,a)+c$, solve
\[
 U^2+cU+a^n=0.
\]
For a nonzero solution $U$ and $u=z_n^k U^{1/n}$, equation~\eqref{eq:dickson}
gives the candidate $x=u+a/u$. A preceding affine centering shift is added back.
The coefficient of $x^{n-2}$ is $-na$, providing an immediate candidate for $a$;
the entire polynomial identity is checked, not inferred from one coefficient.
Even degrees are already eligible for the power-composition reduction.

The dependent term $a/u$ is crucial. Independently selecting an $n$th root of
$U$ and an $n$th root of $a^n/U$ need not produce a pair whose product is $a$.
The correlated form automatically preserves that constraint. This is the
systematic version of the substitution $x=y+1/y$ discussed in the original
question's answers \cite{question}.

\subsection{Generalized reciprocal reduction}
Let $f(x)=\sum_{k=0}^{2q}c_kx^k$ be monic and suppose a nonzero rational $a$
satisfies
\begin{equation}\label{eq:recipcoeff}
 c_{q-j}=a^j c_{q+j}\qquad(1\le j\le q).
\end{equation}
Then
\begin{equation}\label{eq:recip}
 f(x)=x^q g(x+a/x),\qquad
 g(t)=c_q+\sum_{j=1}^q c_{q+j}D_j(t,a).
\end{equation}
Indeed division by $x^q$ pairs the $j$th positive and negative powers, and
$D_j(x+a/x,a)=x^j+a^jx^{-j}$ by~\eqref{eq:dickson}. Thus a degree-$2q$ equation
reduces to a degree-$q$ equation followed by
\[
 x=\frac{t\pm\sqrt{t^2-4a}}2.
\]
The nonzero constant coefficient excludes the lost value $x=0$. The first
nonzero upper coefficient in~\eqref{eq:recipcoeff} yields finitely many rational
candidates for $a$; all coefficient relations are then checked exactly.
The familiar reciprocal and anti-reciprocal-looking cases are only special
instances of this identity.

\subsection{Pair-sum resolvents}
Let $r_1,\ldots,r_n$ be the distinct roots of monic $f$. Define
\[
 R_2(y)=\prod_{i<j}(y-r_i-r_j)\in\Q[y].
\]
A direct resultant identity is
\begin{equation}\label{eq:pairres}
 \Res_x\bigl(f(x),f(y-x)\bigr)
 =\prod_{i,j}(y-r_i-r_j)
 =2^n f(y/2)R_2(y)^2.
\end{equation}
This gives an exact implementation: divide out the diagonal factor, factor over
$\Q$, and halve the exponents. Repeated pair sums remain repeated in $R_2$;
taking only a square-free part would be wrong. The release checks the degree
$n(n-1)/2$ and the exact division. The Python tests verify the full identity
for the motivating sextic.

If a low-degree factor of $R_2$ gives a radical $\beta$, factor $f$ over
$\Q(\beta)$ and solve the smaller factors. The discovered $\beta$ is an extension
candidate; no claim is made that every such factor helps the selected root.
The Python pair path restricts lifts to degrees at most three, and the Wolfram
pair path uses the same restriction, to avoid very large quartic expressions.
User-supplied extensions allow degree-four lifts. Pair resolvents are attempted
only up to a configured original degree (default eight).

General subset sums or elementary symmetric functions would extend this search,
but the degree $\binom nk$ and collisions can become expensive. They are not
necessary for the complete normal-field backend and are not implemented here.

\subsection{Supplied radical extensions}
A caller may supply a known radical element generating a useful field, or in
Wolfram a list of such generators. The package factors over that extension,
lifts small factors, and verifies the target. For Python, a primitive element
such as $\sqrt3+\sqrt5$ is convenient. This does not assert that \emph{every}
sum of suggested generators is primitive; a true compositum can instead be
constructed explicitly before choosing a generator.

The extension must itself pass the radical grammar. Supplying an unsolved
\texttt{Root} as a coefficient field would otherwise merely move the original
problem into the coefficients and falsely present that as a radical solution.

\section{Worked examples}
\subsection{The quintic from the question}
The polynomial is
\[
 f_5(x)=5x^5-25x^3+25x+6
       =5\bigl(D_5(x,1)+6/5\bigr).
\]
Let
\begin{equation}\label{eq:q5}
 u=\left(\frac{-3+4i}{5}\right)^{1/5}.
\end{equation}
Then $u^5+u^{-5}=-6/5$, hence all five candidates
\[
 x_k=z_5^k u+(z_5^k u)^{-1},\qquad 0\le k<5,
\]
satisfy $f_5(x)=0$. Because $|u|=1$, each is real. Put
$\theta=\Arg((-3+4i)/5)/5\in(\pi/10,\pi/5)$. The five angles are
$\theta+2k\pi/5$; elementary comparison of their cosines shows that $k=0$
gives the largest real root. Consequently the target
\begin{lstlisting}
Root[6 + 25 # - 25 #^3 + 5 #^5 &, 5]
\end{lstlisting}
is represented by
\begin{equation}\boxed{u+u^{-1}}.\label{eq:q5answer}\end{equation}
The trigonometric observation is only used here to explain root order; the
package output~\eqref{eq:q5answer} contains no trigonometric expression. The
Python suite independently verifies all five selected roots, not just this one.

\subsection{The sextic becomes a cubic and a quadratic}
For
\[
 f_6(x)=x^6+x^4-x^3-x^2-1
\]
we have the exact identity
\begin{equation}\label{eq:sextic}
 f_6(x)=x^3\left[(x-1/x)^3+4(x-1/x)-1\right].
\end{equation}
Thus take $t=x-1/x$, solve $t^3+4t-1=0$, then solve $x^2-tx-1=0$.
A branch-correlated real Cardano expression is
\begin{equation}\label{eq:cardano}
 C=\left(\frac12+\frac{\sqrt{849}}{18}\right)^{1/3},
 \qquad t=C-\frac{4}{3C}.
\end{equation}
Here $C>0$. Writing $v=-4/(3C)$ gives $Cv=-4/3$ and
$C^3+v^3=1$, which proves $t^3+4t=1$. Its derivative $3t^2+4$ is positive,
so this is the cubic's unique real root. The two real sextic roots are
\[
 \frac{t-\sqrt{t^2+4}}2<0<\frac{t+\sqrt{t^2+4}}2.
\]
The target \code{Root[..., 2]} for this sextic is therefore
\begin{equation}\boxed{\frac{t+\sqrt{t^2+4}}2}.\label{eq:sexticanswer}\end{equation}
The package may print a different but exactly equivalent cubic radical form.
The test criterion is equality, not text identity with~\eqref{eq:cardano}.

The same structure appears in the pair-sum resolvent:
\begin{align*}
 R_2(y)&=(y^3+4y-1)h(y),\\
 h(y)&=y^{12}-y^9+8y^8-2y^7-y^6+4y^5\\
     &\hspace{1.5em}+18y^4+9y^3+2y^2-1.
\end{align*}
Each paired quadratic has root sum $t$. This explains why a low-degree
extension search can discover the cubic even when a direct sextic radical
conversion does not.

\subsection{A degree-eight notebook example}
One explicit fixture from the notebook is
\begin{align*}
 f_8(x)={}&x^8+8x^7+20x^6+8x^5-36x^4\\
         &-48x^3-15x^2+2x+1.
\end{align*}
Put $s=2x+2$. Its equation becomes
\[
 s^8-32s^6+224s^4-448s^2+256=0.
\]
This identity can also be obtained by eliminating $a,b$ from
$a^2=3$, $(b^2-5)^2=20$, and $s=a+b$; the coefficient normalization is
verified in \code{examples/identities.py} rather than assumed from numerical matching.
Its smallest real root has the radical form
\begin{equation}\label{eq:degree8}
 \frac{-2-\sqrt3-\sqrt{5+2\sqrt5}}2.
\end{equation}
The eight conjugates result from independent choices of the signs of
$\sqrt3$, $\sqrt5$, and the outer square root. They are distinct, and
all radicands $5\pm2\sqrt5$ are positive. Choosing both outer contributions
negative and the larger radicand gives the smallest value. The Python regression
uses the hint $\sqrt3+\sqrt5$, factors exactly over that field, and checks both
its selected-root identity and equality with~\eqref{eq:degree8}. This is a
\emph{supplied-extension} success, not a claim that the generic normal-field
route was benchmarked on this degree-eight field.

\subsection{Forcing the general method}
A useful control experiment is $f=x^3-2$. Its splitting field has degree six.
After adjoining the roots of unity specified by $m=6$, the relative group over
$\Q(z_6)$ has order three. The general backend constructs a single prime-index
Kummer step. For the real root it returns an expression equal to $2^{1/3}$;
a separate test targets a nonreal conjugate.

For $x^4-2$, the generic test builds the degree-eight splitting field. Since
$m=2$, the cyclotomic base is just $\Q$, and the observed chain has three
quadratic steps. The final expression is exact but much less attractive than
$-2^{1/4}$ for the smallest real root. This is expected: the purpose of this
test is to exercise the normal-field construction, not to compete with the
power-composition fast path.

\section{Using the software}
\subsection{Wolfram Language}
From the extracted archive directory:
\begin{lstlisting}
Get["wolfram/RadicalRoot.wl"];
a = Root[6 + 25 # - 25 #^3 + 5 #^5 &, 5];
e = RootToRadicals[a, Method -> "Fast"];
VerifyRadical[a, e]
\end{lstlisting}
The expected verification result is \code{True}; this usage is supplied for
local Wolfram execution, not reported as executed here. To retain construction
data and force the general backend:
\begin{lstlisting}
r = RadicalSolve[Root[#^3 - 2 &, 1],
  Method -> "Galois", TimeConstraint -> 300];
If[AssociationQ[r], r["Certificate"], r]
\end{lstlisting}
The three method names are strings: \code{"Automatic"}, \code{"Fast"}, and
\code{"Galois"}. The default field-degree and relative-group-order limits are
48 and 96. The pair-resolvent degree bound is eight. Disable mathematical size
bounds explicitly with
\begin{lstlisting}
"MaxFieldDegree" -> Infinity,
"MaxGroupOrder" -> Infinity
\end{lstlisting}
and the whole-call time limit with \code{TimeConstraint -> Infinity}. Removing
these limits can cause very long computation and memory exhaustion. A primitive
field must be constructed before its degree can be measured, so a degree bound
alone cannot bound the cost of that initial operation.

\code{RadicalSolve} returns an association on success or a \code{Failure}.
\code{RootToRadicals} extracts the expression on success and propagates a
failure otherwise. It never silently returns the input root as if that were a
successful radical conversion. To run the supplied tests:
\begin{lstlisting}
TestReport["tests/RadicalRoot.wlt"]
\end{lstlisting}

\subsection{Python API and command line}
Install the included Python project in an appropriate environment:
\begin{lstlisting}
python -m pip install ./python
python -m pip install pytest
PYTHONPATH=python python -m pytest tests/test_radicalroot.py -v
\end{lstlisting}
On Windows, installation makes the \code{PYTHONPATH} prefix unnecessary. The
Python API is
\begin{lstlisting}
import sympy as s
from radicalroot import radicalize, verify_result
x = s.Symbol("x")
r = radicalize(5*x**5 - 25*x**3 + 25*x + 6,
               index=4, method="fast")
assert r.success and verify_result(r)
print(r.expression)
\end{lstlisting}
\textbf{Index warning:} Python indices are SymPy's zero-based
\code{CRootOf} indices. Wolfram indices are one-based and their nonreal-root
ordering must not be assumed to agree with SymPy's. The two real examples here
use the corresponding shift only because their relevant real-root orders are
known. No general index-only conversion bridge is supplied.

The direct Python API has no hard wall-clock timeout. Use the command-line
wrapper, which runs the calculation in a separate process:
\begin{lstlisting}
python -m radicalroot --coefficients "5,0,-25,0,25,6" \
  --index 4 --method fast --timeout 30
python -m radicalroot --coefficients "1,0,0,-2" \
  --index 0 --method galois --timeout 120
\end{lstlisting}
Coefficients are descending rational numbers separated by commas. They are
parsed with \code{fractions.Fraction}, not Python \code{eval}. The output is JSON.
A zero command-line size or time limit disables that limit; in the Python API,
\code{None} disables a degree or group-order bound. A timeout terminates the
worker and reports a resource limit, not nonsolvability.

\subsection{Result categories}
\begin{center}
\begin{tabular}{@{}p{.23\textwidth}p{.68\textwidth}@{}}
\toprule
Python status & Meaning\\
\midrule
\code{success} & Allowed radical expression constructed with exact equality
checks; includes method and audit data.\\[3pt]
\code{not\_solvable} & An exact nonsolvable Galois group or nontrivial stable
derived subgroup has been computed.\\[3pt]
\code{not\_found} & Bounded fast paths have been exhausted. No conclusion about
existence.\\[3pt]
\code{resource\_limit} & A configured degree, order, or worker-time bound was
reached. No conclusion about existence.\\[3pt]
\code{backend\_failure} & Required exact CAS operation or internal consistency
check did not complete successfully.\\[3pt]
\code{unsupported\_input} & Input or options are outside the stated contract.\\
\bottomrule
\end{tabular}
\end{center}
Wolfram uses analogous capitalized tags. It also distinguishes a final
\code{VerificationFailure}. For degrees up to six, Python can first
use SymPy's documented Galois-group routine as an exact negative shortcut
\cite{sympyfields}. A solvable group from that shortcut is not itself treated
as a radical expression; construction still proceeds. There is no degree-six
restriction on the separately implemented explicit normal-field route.

\section{Validation and its limits}\label{sec:validation}
The supplied validation records identify the exact software environment and
commands. Tests exercise radical grammar rejection, Dickson identities for
$n=2,\ldots,9$, all five quintic conjugates, both real sextic roots, affine
power families, factor selection on reducible input, exact pair-resolvent
multiplicities, group-chain construction, nonsolvability, invalid input,
resource limits, wrong-conjugate rejection, and the notebook degree-eight
extension fixture.

The forced general-path fixtures include a quadratic, the real and one nonreal
root of $x^3-2$, the cyclotomic polynomial $x^4+x^3+x^2+x+1$, and $x^4-2$.
Every successful fixture invokes the independent final equality verifier, not
just a polynomial-residual test. The finite-group tests include cyclic,
dihedral, symmetric and alternating solvable groups, and rejection of $A_5$.
The polynomial $x^5-x-1$ supplies an exact nonsolvable negative test using
SymPy's Galois-group computation.

\begin{center}
\begin{tabular}{@{}ll@{}}
\toprule
Recorded environment/result & Value\\
\midrule
Python & 3.13.5\\
SymPy & 1.14.0\\
pytest & 9.0.2\\
Final Python regression run & 48 passed, 0 failed\\
Aggregate elapsed test time & 40.22 seconds\\
Slowest fixture & General $x^4-2$ path, 22.96 seconds\\
Wolfram runtime tests & Not executed: evaluation service returned HTTP 404\\
\bottomrule
\end{tabular}
\end{center}
These timings describe this single preparation environment and run. They are
neither cross-platform benchmarks nor worst-case bounds. The records are
\path{validation/pytest.xml} (JUnit) and
\path{validation/pytest-output.txt} (console output).


Earlier development runs exposed two ordinary validation defects: an uncaught
constant-polynomial construction exception and a test comparing identical
polynomials over different declared coefficient domains. Both were corrected
before the final recorded test run. A preliminary aggregate run also exceeded
its external time budget; the final run, not that interrupted run, is the basis
of the reported pass count. This history should not be confused with an
all-input performance guarantee.

The Wolfram code has been reviewed and supplied with local tests, but no Wolfram
runtime pass count is claimed. Neither implementation has been exhaustively
validated over all solvable transitive groups, arbitrary high degrees, or the
entire notebook. The current Python backend can spend substantial time simply
constructing a primitive element from several high-degree \code{CRootOf}
objects. Failures and resource limits remain first-class outcomes.

\section{Cost, optimization, and what is not implemented}
The expensive parameter is often the splitting-field degree, not the target's
minimal-polynomial degree $n$. In general $d_L\le n!$, and
\[
 D=[L\Q(z_m):\Q]\le d_L\varphi(m).
\]
This bound does not predict practical running times: coefficient heights,
primitive-element choices, algebraic-number membership and expression expansion
can dominate. The explicit automorphism table has $|H|^2$ entries. At most
$D$ trace candidates are examined per chain step; each involves exact arithmetic
in a $D$-dimensional rational vector space. The coordinate solves also operate
in that ambient space. These are finite but not lightweight computations.

Several improvements are natural but are \emph{not} implemented claims of this
release. A Stauduhar-style permutation Galois-group backend can avoid constructing
the full splitting field merely to diagnose nonsolvability; such algorithms
are developed by Fieker and Kl\"uners \cite{fieker}. Relative field towers and
subgroup-specific invariants could avoid repeated work in a single very large
primitive field. Primitive-element reduction and coefficient-height control
would make exact arithmetic more practical. A subgroup-chain cost function
could prefer small radicands rather than the first available prime quotient.
Finally, interval certificates and a serialized quotient-field witness could
support a small independent verifier.

A target-specific solver could stop before constructing all of $M$ when the
target already lies in a constructed fixed field. The present implementation
constructs the full chain for simplicity. It likewise does not minimize the
number of radicals, radical nesting depth, algebraic coefficient height, or
printed expression size. A shorter answer may require a different extension,
a nonobvious change of primitive generator, or a denesting stage after
construction. Those optimizations should preserve the exact branch invariants
rather than applying unrestricted power identities.

\section{Conclusion}
The notebook's extension searches are best understood as useful resolvent
heuristics, not a complete criterion for radical solvability. The systematic
replacement has a clear mathematical core: compute the normal field, reduce to
prime-cyclic fixed-field steps over a sufficient cyclotomic base, construct each
radical generator by a finite trace--Fourier search, and carry exact embedding
information through every principal-power choice.

This gives an implemented general reference algorithm rather than a degree-four
wrapper or an enumeration of increasingly complicated guessed expressions. The
structural fast paths solve the motivating examples much more economically.
The release deliberately separates three claims: a complete mathematical
construction under exact backend contracts, executable bounded reference
software, and the finite set of tests actually run. That separation is essential
for making the package useful without presenting a timeout or an untested
backend as a theorem about radicals.

\appendix
\section{Reproduction and source map}
\begin{center}
\begin{tabular}{@{}p{.43\textwidth}p{.48\textwidth}@{}}
\toprule
Path & Purpose\\
\midrule
\code{wolfram/RadicalRoot.wl} & Native exact Wolfram implementation.\\
\code{python/radicalroot/core.py} & Structural and Galois--Kummer algorithms.\\
\code{python/radicalroot/cli.py} & Rational-input CLI with worker timeout.\\
\code{tests/test\_radicalroot.py} & Executed Python regression tests.\\
\code{tests/RadicalRoot.wlt} & Wolfram tests for local execution.\\
\code{examples/basic.py}, \code{basic.wl} & Entry-point examples.\\
\code{validation/} & Actual test logs and status notes.\\
\code{provenance/} & Notebook extraction and input hashes.\\
\code{tools/extract\_notebook.py} & Non-evaluating inspection extractor.\\
\bottomrule
\end{tabular}
\end{center}
Compile the article by running \code{make pdf} or two passes of
\code{pdflatex radicalroot.tex} in \code{article/}. The small
\code{validation\_summary.tex} file is part of the article source and must remain
beside it. No font files, external images, or external bibliography databases
are needed. The original input archive is identified by hash but is not
repackaged as newly licensed source.

\section{A schematic algorithm}
\begin{lstlisting}
RADICALIZE(alpha):
  f := minimal polynomial of the selected embedded alpha
  try structural radical candidates, accepting exact equality only
  L := the exact splitting field of f in the fixed complex embedding
  m := product of distinct primes dividing [L:Q]
  M := compositum(L, Q(zeta_m)); choose an integral primitive theta
  A := all exact automorphisms of M, by splitting minpoly(theta) in M
  H := automorphisms in A fixing zeta_m
  if H has a nontrivial stable derived subgroup:
      return NOT_SOLVABLE
  chain := prime-cyclic subnormal chain of H
  B := rational power basis of Q(zeta_m)
  expressions := matching symbolic basis
  for S > N in chain, p = [S:N], sigma outside N:
      find first nonzero rho among the D trace-Fourier basis candidates
      verify N-fixedness, sigma(rho) = zeta_p*rho, S-fixedness of rho^p
      express rho^p in B by exact rational linear algebra
      determine rho / principal_root(rho^p,p) exactly in {zeta_p^b}
      define the corresponding branch-correct radical
      enlarge B and its expression basis by rho^j, 0 <= j < p
  express alpha in the final basis and substitute radical definitions
  verify allowed syntax and exact embedded value
  return expression and audit data
\end{lstlisting}
The actual wrappers surround this scheme with explicit input checks and resource
limits. No ``try until lucky'' step remains in the complete mathematical core.

\begin{thebibliography}{9}
\bibitem{question}
V. Reshetnikov, \emph{Why is ToRadicals not able to handle all cases? Is there a
workaround?} Mathematica Stack Exchange, question 34011, including its discussion
and answers. \url{https://mathematica.stackexchange.com/q/34011}.
Accessed 8 September 2026. The two motivating polynomials are from this question;
worked derivations here are independently presented.

\bibitem{milne}
J. S. Milne, \emph{Fields and Galois Theory}, version 5.10, September 2022.
In particular, independence of homomorphisms, Proposition 5.27, and Theorem 5.34.
\url{https://www.jmilne.org/math/CourseNotes/FT.pdf}.

\bibitem{fieker}
C. Fieker and J. Kl\"uners, \emph{Computation of Galois groups of rational
polynomials}, LMS Journal of Computation and Mathematics \textbf{17}(1)
(2014), 141--158. Preprint \url{https://arxiv.org/abs/1211.3588}.

\bibitem{toradicals}
Wolfram Research, \emph{ToRadicals}, Wolfram Language Documentation.
\url{https://reference.wolfram.com/language/ref/ToRadicals.html}.
Accessed 8 September 2026.

\bibitem{tonumberfield}
Wolfram Research, \emph{ToNumberField}, Wolfram Language Documentation.
\url{https://reference.wolfram.com/language/ref/ToNumberField.html}.
Accessed 8 September 2026.

\bibitem{wlexact}
Wolfram Research, \emph{RootReduce}, \emph{FactorList}, and
\emph{AlgebraicNumberPolynomial}, Wolfram Language Documentation.
\url{https://reference.wolfram.com/language/ref/RootReduce.html};
\url{https://reference.wolfram.com/language/ref/FactorList.html};
\url{https://reference.wolfram.com/language/ref/AlgebraicNumberPolynomial.html}.
Accessed 8 September 2026.

\bibitem{sympyfields}
SymPy developers, \emph{Number Fields}, SymPy 1.14.0 documentation.
\url{https://docs.sympy.org/latest/modules/polys/numberfields.html}.
The documented Galois-group shortcut has degree bound six; the number-field
operations support the separate general construction in this release.
Accessed 8 September 2026.

\bibitem{sympydomains}
SymPy developers, \emph{Domains Reference}, SymPy 1.14.0 documentation.
\url{https://docs.sympy.org/latest/modules/polys/domainsref.html}.
Algebraic fields and exact quotient-field elements.
Accessed 8 September 2026.
\end{thebibliography}
\end{document}
